\newpage
\phantomsection
\label{prf:euclid-i-1}
\LRAProofFor{thm:euclid-i-1}

\begin{remark*}[Return]
\hyperref[thm:euclid-i-1]{Return to Theorem}
\end{remark*}

\begin{theorem*}[Euclid I.1: Constructing an equilateral triangle]
Given a finite straight line segment $AB$, there exists an equilateral
triangle $ABC$ constructed on $AB$.
\end{theorem*}

\begin{proof}[Professional Standard Proof]
\LRAProofBodyStart
Let $AB$ be the given finite straight line. Describe the circle with center
$A$ and radius $AB$, and describe the circle with center $B$ and radius $BA$.
Let $C$ be an intersection point of the two circles, and draw the straight
lines $CA$ and $CB$.

Since $C$ lies on the circle with center $A$, $AC=AB$. Since $C$ lies on the
circle with center $B$, $BC=BA$. Thus $AC$ and $BC$ are each equal to $AB$,
so by Common Notion I they are equal to one another. Therefore $AB$, $BC$,
and $CA$ are equal, and the triangle $ABC$ is equilateral.
\end{proof}

\begin{proof}[Detailed Learning Proof]
\LRAProofBodyStart
\begin{center}
\begin{tikzpicture}[scale=1.35, line cap=round, line join=round]
  \definecolor{euclidblue}{RGB}{30,110,190}
  \tikzset{
    equaltick/.style={
      postaction={decorate},
      decoration={markings, mark=at position 0.5 with {
        \draw[line width=0.6pt] (0,-2.4pt)--(0,2.4pt);
      }}
    }
  }

  \coordinate (A) at (0,0);
  \coordinate (B) at (2.2,0);
  \pgfmathsetmacro{\R}{2.2}
  \coordinate (C) at (1.1,{sqrt(3)/2*\R});

  \draw[gray!55, thin] (A) circle (\R);
  \draw[gray!55, thin] (B) circle (\R);

  \draw[very thick, equaltick] (A) -- (B);
  \draw[very thick, euclidblue, equaltick] (A) -- (C);
  \draw[very thick, euclidblue, equaltick] (B) -- (C);

  \foreach \p/\pos in {A/below left,B/below right,C/above}{
    \fill (\p) circle (0.025);
    \node[\pos] at (\p) {$\p$};
  }

  \node[gray!70, font=\footnotesize] at ($(A)+(-1.00,-1.10)$) {center $A$};
  \node[gray!70, font=\footnotesize] at ($(B)+(1.00,-1.10)$) {center $B$};
  \node[euclidblue, font=\footnotesize] at ($(A)!0.5!(C)+(-0.24,0.06)$) {$=AB$};
  \node[euclidblue, font=\footnotesize] at ($(B)!0.5!(C)+(0.28,0.06)$) {$=AB$};
\end{tikzpicture}
\end{center}


The construction uses exactly the permitted tools available at this point.
Postulate I allows the given segment $AB$ to be treated as a drawable straight
line segment. Postulate III allows a circle with any center and distance, so we
may draw one circle centered at $A$ with distance $AB$ and another centered at
$B$ with distance $BA$.

Choose an intersection point $C$ of the two circles. The segment $AC$ is a
radius of the circle centered at $A$, so it equals the radius $AB$. The segment
$BC$ is a radius of the circle centered at $B$, so it equals the radius $BA$,
which is the same segment as $AB$. Hence all three sides of $\triangle ABC$
are equal, and by the definition of an equilateral triangle, $\triangle ABC$ is
equilateral.
\end{proof}

\begin{remark*}[Proof structure]
Construct two circles from the endpoints of the given segment, draw the two
radii to an intersection point, and use equality through the common segment
$AB$ to identify the three sides.
\end{remark*}

\begin{remark*}[Euclid's intersection gap]
The proof assumes that the two circles meet. This is the classical lacuna in
Euclid I.1: the postulates permit drawing the circles, but Euclid's stated
axioms do not by themselves include a circle--circle intersection principle.
Later axiomatizations add continuity or incidence assumptions that justify this
step.
\end{remark*}

\begin{dependencies}
\begin{itemize}
  \item \hyperref[def:euclid-point]{Point}
  \item \hyperref[def:euclid-straight-line]{Straight line}
  \item \hyperref[def:euclid-circle]{Circle}
  \item \hyperref[def:euclid-triangle-side-classes]{Equilateral triangle}
  \item \hyperref[ax:euclid-postulate-draw-straight-line]{Euclid Postulate I}
  \item \hyperref[ax:euclid-postulate-describe-circle]{Euclid Postulate III}
  \item \hyperref[ax:euclid-common-notion-things-equal-to-same]{Euclid Common Notion I}
\end{itemize}
\end{dependencies}

\clearpage
